\section{差分と項数の削減}\label{節：差分と項数の削減}
\subsection{差を取ると項数が減る理由}

前節では$a_{n+1}-2a_n$を一つの数列として追った.
まず倍率が1の場合,つまり隣り合う二項の差を考えよう.
初期値$a_1=1,a_2=2,a_3=6$の下で
\[
 a_{n+3}=6a_{n+2}-11a_{n+1}+6a_n
\]
を左辺へまとめると
\[
 \begin{aligned}
 &(a_{n+3}-a_{n+2})-5(a_{n+2}-a_{n+1})\\
 &\qquad+6(a_{n+1}-a_n)=0.
 \end{aligned}
\]
従って$b_n=a_{n+1}-a_n$と置けば
\[
 b_{n+2}-5b_{n+1}+6b_n=0
\]
という隣接3項間漸化式になる.
元の隣接4項間漸化式より,必要な初期値が一つ少ない.
ただし差だけでは元の数列を一意に決められない.
$b_n$を解いた後に
\[
 a_n=a_1+\sum_{j=1}^{n-1}b_j
\]
と戻すため,残る初期値$a_1$も使う.
$n=1$では和を0とする.


初期値から$b_1=1,b_2=4$となる.
前節の例題と同じ式・初期値なので,
\[
 b_n=-2^{n-1}+2\cdot3^{n-1}.
\]
和を取ると
\[
 \begin{aligned}
 a_n&=1-(2^{n-1}-1)+(3^{n-1}-1)\\
 &=1-2^{n-1}+3^{n-1}.
 \end{aligned}
\]
差分を取る前の初期値$a_1$を用いて,定数の成分も取り戻せた.

なぜこの例では差が有効なのだろうか.
対応する多項式は
\[
 X^3-6X^2+11X-6=(X-1)(X-2)(X-3)
\]
であり,差を取ることが因子$X-1$に対応する.

\subsection{項を一つ進める操作}

数列の添字を一つ進める操作を$E$と書く.
すなわち$(Ea)_n=a_{n+1}$であり,
\[
 ((E-r)a)_n=a_{n+1}-ra_n
\]
と定める. ここで$E$は数ではなく,数列に作用する記号である.
二度進める操作は$E^2$,何もしない操作は$I$と書く.
表記を簡単にするため,$E-rI$を$E-r$と略している.

例えば前の隣接4項間漸化式は
\[
 (E^3-6E^2+11E-6)a=0
\]
であり,多項式と同じように因数分解して
\[
 (E-2)(E-3)\{(E-1)a\}=0
\]
と書ける. 右辺の0は全項が0の数列を表す.
$b=(E-1)a$と置けば,残った因子が$b$の漸化式を与える.
各因子は同じ$E$の多項式なので,作用させる順序を交換できる.

一般に,特性多項式$P(X)$が
\[
 P(X)=(X-r)Q(X)
\]
と分かれるなら,$b_n=a_{n+1}-ra_n$によって
\[
 P(E)a=0\quad\Longleftrightarrow\quad Q(E)b=0
\]
となる. 特に通常の差分$\Delta a_n=a_{n+1}-a_n$が使えるのは,$P(1)=0$の場合である.
根が1でない式に,通常の差分を取るだけで項数が減るとは限らない.

\subsection{倍率を取り除いてから差を取る}

$r\ne0$とし,$a_n=r^{n-1}c_n$と置けば
\[
 a_{n+1}-ra_n=r^n(c_{n+1}-c_n).
\]
従って$a_{n+1}-ra_n$も,倍率$r$の等比的な変化を取り除いた後の差分である.
異なる根を一つずつ取り除く方法も,重解の場合に差を繰り返す方法も,同じ因子の計算で説明できる.

\begin{bookpoint}
行列では,更新を一定倍率の成分に分ける.
差分では,その倍率に対応する因子を一つずつ取り除く.
項数が減る代わりに,元へ戻すときに初期値を一つ使う.
\end{bookpoint}
